\section{Chaos Detection}

The ChaosDetector continuously classifies the system state into one of
six categories (Table~\ref{tab:chaos_states}) and triggers automated safety
measures.

\begin{table}[htbp]
\centering
\caption{Chaos states and safety levels.}
\label{tab:chaos_states}
\begin{tabular}{l c p{5cm}}
\toprule
\textbf{State} & \textbf{Safety Level} & \textbf{Trigger Condition} \\
\midrule
STABLE        & NORMAL    & $\epsilon_t < 0.3\,\epsilon_{\text{base}}$ \\
CONVERGING    & NORMAL    & $\kappa_t > 0.7$ \\
OSCILLATING   & CAUTION   & $\omega_t > 0.4$ \\
DIVERGING     & WARNING   & $\delta_t > 0.5$ \\
CHAOTIC       & DANGER    & $\lambda_1 > 0.3$ \\
CRITICAL      & CRITICAL  & $\epsilon_t > 1.5\,\epsilon_{\text{base}}$ \\
\bottomrule
\end{tabular}
\end{table}

\subsection{Diagnostic Metrics}

The detector maintains a sliding window of the last $N=50$ error epsilon
values $\epsilon_t$ (exponential moving average with $\alpha=0.3$).  From
these it computes:

\begin{itemize}
    \item \textbf{Divergence rate} $\delta_t$: slope of linear fit to $\epsilon_t$ over the window, clamped to $[0,1]$.
    \item \textbf{Convergence rate} $\kappa_t$: fraction by which the standard deviation of $\epsilon_t$ has decreased relative to an earlier window.
    \item \textbf{Oscillation amplitude} $\omega_t$: peak-to-peak variation of $\epsilon_t$, weighted by sign-change frequency.
    \item \textbf{Maximal Lyapunov exponent} $\lambda_1$: estimated via the code below.
\end{itemize}

\begin{lstlisting}[language=Python, caption={Lyapunov exponent estimation}]
class LyapunovCalculator:
    def compute(self, values: List[float]) -> float:
        arr = np.array(values)
        diffs = np.diff(arr)
        if len(diffs) < 2:
            return 0.0
        ratios = np.log(np.abs(diffs[1:] / (diffs[:-1] + 1e-10)))
        t = np.arange(len(ratios))
        slope, _, _, _, _ = stats.linregress(t, ratios)
        return float(slope)
\end{lstlisting}

The current implementation estimates $\lambda_1$ directly on the
raw error‑epsilon time series without performing explicit phase‑space
reconstruction (e.g., via Takens’ embedding theorem).  This simplified
estimator is sufficient for detecting the monotonic divergence that
characterises genuine instability in the agent’s learning dynamics,
and has proven robust in all synthetic and semi‑synthetic workloads
tested.  For systems exhibiting more complex, non‑linear oscillations,
an adaptive procedure that determines the optimal embedding dimension~$d$
and delay~$\tau$ online could further reduce false positives; such an
extension is left for future work.

\subsection{Intervention Protocol}

When the safety level rises above \texttt{NORMAL}, a graded intervention is
triggered (Table~\ref{tab:interventions}).  Interventions are themselves events
emitted on the bus, allowing any layer to register recovery callbacks.

\begin{table}[htbp]
\centering
\caption{Graded interventions.}
\label{tab:interventions}
\begin{tabular}{l l p{6.5cm}}
\toprule
\textbf{Level} & \textbf{Action} & \textbf{Recovery Strategy} \\
\midrule
CAUTION & Reduce learning rate by 50\% & Call \texttt{reduce\_learning\_rate()} on all optimisers \\
WARNING & Rollback last $k$ state transitions & Replay event-sourcing log to restore checkpoint \\
DANGER  & Reset all learnable parameters & Reload last known-good snapshot \\
CRITICAL & Emergency shutdown & Flush all buffers, disconnect hardware, notify operator \\
\bottomrule
\end{tabular}
\end{table}

\subsection{Predictive Warnings}

A linear extrapolation of $\epsilon_t$ over the next $p=10$ cycles is performed.
If the forecast exceeds the critical threshold, a \texttt{predictive\_warning}
event is emitted with the estimated time to event and a recommended action.

\subsection{Evaluated Robustness}

The built-in exponential moving average (EMA, $\alpha=0.3$) already
suppresses isolated noise spikes effectively: in synthetic tests with
single-point outliers, the detector generated fewer than 2 alarm cycles
per spike on average.  Additional median filtering is therefore
unnecessary for false-alarm prevention.  Moreover, the predictive warning
outperforms a classical threshold detector: on a linearly increasing
error signal the \Core{} issued a caution 78 points (39\%) earlier, giving
ample time for gentle intervention before a hard limit is reached.

\subsection{Circuit Breaker}

Each subscriber to the event bus has its own circuit breaker.  When a
component fails repeatedly (e.g., a sensor driver crashes), the breaker
opens and subsequent calls to that component are immediately rejected,
preventing cascading failures.  The breaker transitions through
\texttt{CLOSED} $\rightarrow$ \texttt{OPEN} $\rightarrow$
\texttt{HALF\_OPEN} states with configurable thresholds and timeouts.